\documentclass[10pt,a4paper]{amsart}
\usepackage{amsmath,amssymb,amsthm,graphicx}
% Retained preamble for the Stacks Project chapter snapshot in this directory.
% The chapter's own preamble.tex is a copy of the shared Stacks Project preamble; the block
% below is sliced MECHANICALLY, line for line, from it. Nothing was added, removed or
% redefined; the wrapper class, the xr-hyper/hyperref wiring and the externaldocument list
% are omitted because the assembler supplies the document class.
%
% lines taken from preamble.tex: 3, 5, 9, 11, 12, 156, 157, 158, ... (51 lines total)
%
% package coverage restored from the chapter's own preamble (2026-10-13): xy, xy-2cell

\usepackage[all]{xy}
\xyoption{2cell}
\UseAllTwocells

\newenvironment{reference}{\comment}{\endcomment}
\newenvironment{slogan}{\comment}{\endcomment}
\newenvironment{history}{\comment}{\endcomment}
\theoremstyle{plain}
\newtheorem{theorem}[subsection]{Theorem}
\newtheorem{proposition}[subsection]{Proposition}
\newtheorem{lemma}[subsection]{Lemma}
\theoremstyle{definition}
\newtheorem{definition}[subsection]{Definition}
\newtheorem{example}[subsection]{Example}
\newtheorem{exercise}[subsection]{Exercise}
\newtheorem{situation}[subsection]{Situation}
\theoremstyle{remark}
\newtheorem{remark}[subsection]{Remark}
\newtheorem{remarks}[subsection]{Remarks}
\numberwithin{equation}{subsection}
\def\lim{\mathop{\mathrm{lim}}\nolimits}
\def\colim{\mathop{\mathrm{colim}}\nolimits}
\def\Spec{\mathop{\mathrm{Spec}}}
\def\Hom{\mathop{\mathrm{Hom}}\nolimits}
\def\Ext{\mathop{\mathrm{Ext}}\nolimits}
\def\SheafHom{\mathop{\mathcal{H}\!\mathit{om}}\nolimits}
\def\SheafExt{\mathop{\mathcal{E}\!\mathit{xt}}\nolimits}
\def\Sch{\mathit{Sch}}
\def\Mor{\mathop{\mathrm{Mor}}\nolimits}
\def\Ob{\mathop{\mathrm{Ob}}\nolimits}
\def\Sh{\mathop{\mathit{Sh}}\nolimits}
\def\NL{\mathop{N\!L}\nolimits}
\def\CH{\mathop{\mathrm{CH}}\nolimits}
\def\proetale{{pro\text{-}\acute{e}tale}}
\def\etale{{\acute{e}tale}}
\def\QCoh{\mathit{QCoh}}
\def\Ker{\mathop{\mathrm{Ker}}}
\def\Im{\mathop{\mathrm{Im}}}
\def\Coker{\mathop{\mathrm{Coker}}}
\def\Coim{\mathop{\mathrm{Coim}}}
\DeclareMathSymbol{\boxtimes}{\mathbin}{AMSa}{"02}
\def\QCohstack{\mathcal{QC}\!\mathit{oh}}
\def\Cohstack{\mathcal{C}\!\mathit{oh}}
\def\Spacesstack{\mathcal{S}\!\mathit{paces}}
\def\Quotfunctor{\mathrm{Quot}}
\def\Hilbfunctor{\mathrm{Hilb}}
\def\Curvesstack{\mathcal{C}\!\mathit{urves}}
\def\Polarizedstack{\mathcal{P}\!\mathit{olarized}}
\def\Complexesstack{\mathcal{C}\!\mathit{omplexes}}
\def\Pic{\mathop{\mathrm{Pic}}\nolimits}
\def\Picardstack{\mathcal{P}\!\mathit{ic}}
\def\Picardfunctor{\mathrm{Pic}}
\def\Deformationcategory{\mathcal{D}\!\mathit{ef}}
\usepackage{cleveref}
\crefname{originallemma}{original lemma}{original lemmas}
\Crefname{originallemma}{Original lemma}{Original lemmas}
\makeatletter
\newlabel{algebra-lemma-NAK}{{00DV}{source}{}{sourcepointer-1}{}}
\newlabel{algebra-lemma-NAK@cref}{{[originallemma][0][]00DV}{[1][0][]source}{}{}{}}
\makeatother
\crefname{originallemma}{original lemma}{original lemmas}
\Crefname{originallemma}{Original lemma}{Original lemmas}
\makeatletter
\newlabel{algebra-lemma-complete-modulo-nilpotent}{{0DYC}{source}{}{sourcepointer-2}{}}
\newlabel{algebra-lemma-complete-modulo-nilpotent@cref}{{[originallemma][0][]0DYC}{[1][0][]source}{}{}{}}
\makeatother
\crefname{originallemma}{original lemma}{original lemmas}
\Crefname{originallemma}{Original lemma}{Original lemmas}
\makeatletter
\newlabel{lemma-lift-equivalence-module}{{0DZ2}{source}{}{sourcepointer-3}{}}
\newlabel{lemma-lift-equivalence-module@cref}{{[originallemma][0][]0DZ2}{[1][0][]source}{}{}{}}
\makeatother
\crefname{originallemma}{original lemma}{original lemmas}
\Crefname{originallemma}{Original lemma}{Original lemmas}
\makeatletter
\newlabel{smoothing-proposition-lift-smooth}{{07M8}{source}{}{sourcepointer-4}{}}
\newlabel{smoothing-proposition-lift-smooth@cref}{{[originallemma][0][]07M8}{[1][0][]source}{}{}{}}
\makeatother
\title{Stacks, Examples of Deformations: lemma-lift-equivalence}
\author{The Stacks Project Authors}
\date{stacks-project:examples-defos.tex:local-snapshot-20261012; GNU FDL 1.2 or later}
\begin{document}
\maketitle
\noindent\textit{Editorial note.} Statement span is lines 2292-2300 of sources/stacks-examples-defos/pinned-original-source.tex; proof span is lines 2302-2351, the complete author proof (48 body lines).\par\medskip
\noindent\textit{Editorial note.} Navigation stubs: 4 cited-but-not-included labels are supplied as non-load-bearing pointers with their pinned tags where the snapshot tag index has them.\par\medskip
\noindent\textit{Editorial note.} extra\_packages is empty on purpose: the retained preamble already loads xy with options and a second declaration clashes in the assembler wrapper.\par\medskip
\noindent\textit{Editorial note.} Text is retained verbatim, including the author's own typos and placeholders.\par\medskip
\noindent\textit{Editorial note.} Queue index 6 of queue\_examples-defos.json (pb=2583, stmt 2292-2300, proof 2302-2351).\par\medskip
\noindent\textit{Allowed standard prerequisites (editorial).}\begin{itemize}
\item Explicit prerequisite boundary. The cited inputs are prerequisites; the argument in the retained span is not replaced by any cited result.
\end{itemize}\medskip
\ifdefined\hypertarget\hypertarget{sourcepointer-1}{}\fi
\noindent\textit{Source reference.} Original lemma 00DV: sources/stacks-examples-defos/pinned-original-source.tex, source lines 2351--2351. Cited by the retained statement or proof but not part of the collected spans; supplied only as a navigation pointer. The argument does not depend on its content.\par\medskip
\ifdefined\hypertarget\hypertarget{sourcepointer-2}{}\fi
\noindent\textit{Source reference.} Original lemma 0DYC: sources/stacks-examples-defos/pinned-original-source.tex, source lines 2351--2351. Cited by the retained statement or proof but not part of the collected spans; supplied only as a navigation pointer. The argument does not depend on its content.\par\medskip
\ifdefined\hypertarget\hypertarget{sourcepointer-3}{}\fi
\noindent\textit{Source reference.} Original lemma 0DZ2: sources/stacks-examples-defos/pinned-original-source.tex, source lines 2213--2234. Cited by the retained statement or proof but not part of the collected spans; supplied only as a navigation pointer. The argument does not depend on its content.\par\medskip
\ifdefined\hypertarget\hypertarget{sourcepointer-4}{}\fi
\noindent\textit{Source reference.} Original lemma 07M8: sources/stacks-examples-defos/pinned-original-source.tex, source lines 2351--2351. Cited by the retained statement or proof but not part of the collected spans; supplied only as a navigation pointer. The argument does not depend on its content.\par\medskip
\begin{lemma}
\label{lemma-lift-equivalence}
Let $A' \to A$ be a surjective map of Noetherian rings with nilpotent kernel.
Let $A \to B$ be a finite type flat ring map.
Let $\mathfrak b \subset B$ be an ideal such that
$\Spec(B) \to \Spec(A)$ is syntomic on the complement of $V(\mathfrak b)$.
Then $B$ has a flat lift to $A'$ if and only if the $\mathfrak b$-adic
completion $B^\wedge$ has a flat lift to $A'$.
\end{lemma}

\begin{proof}
Choose an $A$-algebra surjection $P = A[x_1, \ldots, x_n] \to B$.
Let $\mathfrak p \subset P$ be the inverse image of $\mathfrak b$.
Set $P' = A'[x_1, \ldots, x_n]$ and denote $\mathfrak p' \subset P'$
the inverse image of $\mathfrak p$. (Of course $\mathfrak p$
and $\mathfrak p'$ do not designate prime ideals here.)
We will denote $P^\wedge$ and $(P')^\wedge$ the respective completions.

\medskip\noindent
Suppose $A' \to B'$ is a flat lift of $A \to B$, in other words,
$A' \to B'$ is flat and there is an $A$-algebra isomorphism
$B = B' \otimes_{A'} A$. Then we can choose an $A'$-algebra map
$P' \to B'$ lifting the given surjection $P \to B$.
By Nakayama's lemma (Algebra, Lemma \ref{algebra-lemma-NAK})
we find that $B'$ is a quotient of $P'$. In particular, we find
that we can endow $B'$ with an $A'$-flat $P'$-module structure
lifting $B$ as an $A$-flat $P$-module.
Conversely, if we can lift $B$ to a $P'$-module $M'$ flat over $A'$,
then $M'$ is a cyclic module $M' \cong P'/J'$ (using Nakayama again)
and setting $B' = P'/J'$ we find a flat lift of $B$ as an algebra.

\medskip\noindent
Set $C = B^\wedge$ and $\mathfrak c = \mathfrak bC$.
Suppose that $A' \to C'$ is a flat lift of $A \to C$.
Then $C'$ is complete with respect to the inverse image $\mathfrak c'$
of $\mathfrak c$
(Algebra, Lemma \ref{algebra-lemma-complete-modulo-nilpotent}).
We choose an $A'$-algebra map $P' \to C'$ lifting
the $A$-algebra map $P \to C$. These maps pass through
completions to give surjections $P^\wedge \to C$ and $(P')^\wedge \to C'$
(for the second again using Nakayama's lemma).
In particular, we find that we can endow $C'$ with an $A'$-flat
$(P')^\wedge$-module structure lifting $C$ as an $A$-flat $P^\wedge$-module.
Conversely, if we can lift $C$ to a $(P')^\wedge$-module $N'$ flat over $A'$,
then $N'$ is a cyclic module $N' \cong (P')^\wedge/\tilde J$
(using Nakayama again) and setting $C' = (P')^\wedge/\tilde J$
we find a flat lift of $C$ as an algebra.

\medskip\noindent
Observe that $P' \to (P')^\wedge$ is a flat ring map which
induces an isomorphism $P'/\mathfrak p' = (P')^\wedge/\mathfrak p'(P')^\wedge$.
We conclude that our lemma is a consequence of
Lemma \ref{lemma-lift-equivalence-module} provided we can
show that $B_g$ lifts to an $A'$-flat $P'_g$-module for
$g \in \mathfrak p'$. However, the ring map $A \to B_g$ is syntomic
and hence lifts to an $A'$-flat algebra $B'$ by
Smoothing Ring Maps, Proposition \ref{smoothing-proposition-lift-smooth}.
Since $A' \to P'_g$ is smooth, we can lift $P_g \to B_g$
to a surjective map $P'_g \to B'$ as before and we get what we want.
\end{proof}

\end{document}
